\section{Глава «Программирование на Python»: цели и маршрут}

\subsection*{Цели раздела}
\begin{itemize}
\item Перейти от базовых команд \texttt{arm}/\texttt{takeoff}/\texttt{land} к автономным миссиям.
\item Освоить два подхода к движению: по точкам и по скоростям.
\item Научиться использовать телеметрию, камеру и базовые алгоритмы компьютерного зрения.
\item Подготовить учащихся к проектам ArUco-навигации и трекинга человека.
\end{itemize}

\subsection{Карта исходных скриптов}
\begin{table}[H]
\centering
\begin{tabularx}{\textwidth}{>{\raggedright\arraybackslash}p{0.31\textwidth}>{\raggedright\arraybackslash}X}
\toprule
Группа скриптов & Учебная роль \\
\midrule
\texttt{mission\_examples/*.py} & Базовое пилотирование, точки, прямолинейный полёт, окружность \\
\texttt{get\_telemetry/*.py} & Диагностика батареи, дальномера и локальной позиции \\
\texttt{camera\_examples/*.py} & Поток, кадры, калибровка камеры \\
\texttt{aruco\_examples/*.py} & Детекция маркеров, оценка позы, полёт по визуальному ориентиру \\
\texttt{PioneerHumanTracking/*.py} & Продвинутый CV: скелетная модель, жесты, контур управления \\
\texttt{group\_flight\_examples/*.py} & Многодронные миссии и синхронизация потоков \\
\bottomrule
\end{tabularx}
\caption{Покрытие учебных тем по структуре репозитория}
\end{table}

\subsection{Дидактическая лестница}
\begin{enumerate}
\item Безопасный старт: соединение, арминг, взлёт, посадка.
\item Контроль движения: скорость и локальные точки.
\item Фигуры полёта: прямоугольник, треугольник, окружность.
\item Сенсоры и обратная связь.
\item Компьютерное зрение.
\item Интеграционный проект.
\end{enumerate}

\begin{infobox}[Методический акцент]
Последовательность заданий построена по принципу «одна новая идея за урок»: сначала только одна новая команда API, затем комбинация команд в мини-миссию.
\end{infobox}

\subsection{Минимальные требования к входу в раздел}
\begin{itemize}
\item Базовый синтаксис Python: переменные, циклы, функции, \texttt{try/except}.
\item Навыки запуска скрипта и чтения консольного вывода.
\item Понимание системы координат \texttt{x,y,z,yaw}.
\end{itemize}

\begin{taskbox}[Стартовая диагностика]
Составьте таблицу из трёх команд API Pioneer, которые отвечают за запуск миссии, движение и безопасное завершение. Для каждой команды укажите один риск ошибочного применения.
\end{taskbox}
